% ECONOMICS_PROBLEM: {"id": "L41-613", "title": "Merger effects beyond prices", "jel_code": "L41", "source_id": "OP-0393"}
% !TeX program = xelatex
\documentclass[11pt,a4paper]{article}
\usepackage[margin=23mm]{geometry}
\usepackage{fontspec}
\setmainfont{Latin Modern Roman}
\usepackage{amsmath,amssymb,mathtools}
\usepackage{xcolor,enumitem,booktabs,longtable,array,xurl,hyperref}
\definecolor{muted}{HTML}{53636C}
\hypersetup{colorlinks=true,urlcolor=blue,linkcolor=blue}
\setlength{\parindent}{0pt}
\setlength{\parskip}{5pt}
\newcommand{\R}{\mathbb R}
\newcommand{\E}{\mathbb E}
\newcommand{\N}{\mathbb N}
\newcommand{\ind}{\mathbf 1}
\newcommand{\Var}{\operatorname{Var}}
\newcommand{\Cov}{\operatorname{Cov}}
\newcommand{\supp}{\operatorname{supp}}
\DeclareMathOperator*{\argmin}{arg\,min}
\DeclareMathOperator*{\argmax}{arg\,max}
\newcommand{\entrymeta}[3]{{\sffamily\small\color{muted}\textbf{\texttt{#1}}\quad|\quad #2\par #3\par}}
\newcommand{\Problemref}[1]{\texttt{#1}}
\newcommand{\JELref}[2]{\texttt{#2} (JEL \texttt{#1})}
\begin{document}
\section*{Merger effects beyond prices}
\textbf{Problem ID:} \texttt{L41-613}\quad \textbf{JEL:} \texttt{L41}\par
% BEGIN ECONOMICS STATEMENT
\label{problem:OP-0393}
{\sffamily\small\color{muted}Original subject group: Industrial organization and competition\par}
\label{definitions:problem:30007675}
\textbf{Audit classification:} Published research agenda. Quality among surviving products and total available quality are different estimands. The revised index is defined even when no product survives.
\subsection*{Definitions and precise research target}
Fix the population of horizontal mergers in United States grocery-product markets closing during a specified five-year cohort. For merger \(m\), define two regimes \(a\in\{0,1\}\), where \(a=1\) implements the merger and \(a=0\) keeps the firms independent. Let \(Q_m(a,h)\) be an independently audited quality index and \(V_m(a,h)\) the number of distinct product variants available in its premerger geographic market at horizon \(h\) years. Normalize quality on a common scale and fix the variant definition before seeing outcomes. The primary targets are \(\tau_Q(h)=E[Q_m(1,h)-Q_m(0,h)]\) and \(\tau_V(h)=E[V_m(1,h)-V_m(0,h)]\), for \(h\in\{1,3,5\}\); expectation averages over the defined merger cohort. Include entry, exit and product redesign within each regime, instead of holding the original product set fixed. Identify these effects using comparable markets under an explicit conditional counterfactual-trends assumption, or report sensitivity bounds when that assumption is relaxed. Quality sampling must not depend only on surviving products, since discontinuation is itself an outcome. Quantify whether quality and variety responses differ with premerger substitution and cost complementarities. This specifies a longitudinal empirical agenda: it does not assume that every merger lowers quality, and it distinguishes changes in product count from changes in consumer-valued variety.

\textbf{Definitions, assumptions and mathematical qualifications.} Let each market have a fixed geographic frame and a finite upper bound \(J_{\max}\) on product variants. A variant has a preregistered identity rule independent of merger outcomes. Define \(V_m(a,h)=\sum_{j\in\mathcal J_m(a,h)}1\). Let \(z_j(a,h)\in[0,1]\) be independently audited quality. One well-defined index including exit is \(Q_m(a,h)=J_{\max}^{-1}\sum_{j\in\mathcal J_m(a,h)}z_j(a,h)\), with zero for absence; it measures quality provision, rather than the mean quality among survivors. State whether a different index is desired. For a fixed cohort use specified weights \(\nu_m\), \(\sum_m\nu_m=1\), in both causal contrasts. A conditional counterfactual-trends assumption requires equality of mean no-merger changes between merger and comparison markets given baseline covariates and adequate overlap.
\subsection*{Variants and assumption changes}
\begin{enumerate}
\item \textbf{Editorial, intensive margin.} Restrict to a predeclared set of products that exist under both regimes, and estimate quality effects for that principal stratum; its membership is latent and needs additional assumptions.
\item \textbf{Editorial, valued variety.} Replace product count by a demand-based equivalent-variation index under specified quasi-linear preferences, including new and discontinued products; this requires identified preferences.
\end{enumerate}
\subsection*{References and source audit}
\textbf{Checked source.} 24 August 2024 author draft, Section 6, printed p. 37, sixth research direction. Full primary text retrieved. \url{https://web.stanford.edu/~ayurukog/trendsincompetition.pdf}.
\begin{enumerate}\raggedright
\item\hypertarget{definitions:source:30007675:1}{} Carl Shapiro; Ali Yurukoglu. 2024. \emph{Trends in Competition in the United States: What Does the Evidence Show?}. 24 August 2024 author draft, Section 6, printed p. 37, sixth research direction.
\url{https://web.stanford.edu/~ayurukog/trendsincompetition.pdf}.
DOI: \url{https://doi.org/10.3386/w32762}.
\end{enumerate}

\subsection*{Common conventions: Source mathematical conventions}

These conventions apply to each entry unless explicitly replaced.

\textbf{Models and identification.} A primitive model \(m\) specifies all probability laws, feasible choices, information, equilibrium and measurement rules used by its target. The admissible class \(\mathcal M\) is fixed independently of outcomes. If \(P_O(m)\) is its observable probability law and \(\tau(m)\) the target, its identified set at observable law \(P\) is
\[
 \mathcal I_\tau(P)=\{\tau(m):m\in\mathcal M,\ P_O(m)=P\}.
\]
Point identification means that this set is a singleton. An empty set signals incompatibility of data and assumptions. Coordinate bounds need not describe the joint set. Bounds are sharp only if every retained target is attainable or arbitrarily approachable by compatible models. Finite bounds require explicit restrictions on unobserved quantities; unrestricted confounding, hidden wealth, extrapolation or arbitrary equilibrium selection can yield uninformative sets.

\textbf{Causal interventions.} A potential outcome \(Y_i(p)\) is the outcome under a complete policy regime \(p\), including timing, financing, information and market spillovers. The target population and units are fixed before treatment. With interference, use the complete assignment vector or a justified exposure mapping; individual no-interference notation alone cannot identify an economy-wide intervention. A difference of mean potential outcomes does not require the cross-world joint distribution. Conditioning on post-treatment migration, survival, enrollment or employment changes the estimand and needs separate justification.

\textbf{Statistical statements.} An estimator, confidence set, forecast or test specifies a sampling law, model class and loss/coverage criterion. Uniform \((1-\alpha)\)-coverage means \(\inf_{m\in\mathcal M}\Pr_m\{\tau(m)\in C_n\}\geq1-\alpha\), or an explicitly asymptotic version. The whole parameter space trivially has coverage, so informativeness requires a stated width, risk or power objective. A forecasting improvement is relative to a named benchmark, information set and evaluation population.

\textbf{Equilibrium and optimization.} An equilibrium comprises feasible plans, optimality, beliefs where needed, budget constraints and all market/resource clearing. The primitives or a stated benchmark define each correspondence \(\mathcal E(p;m)\). Nonemptiness is assumed or proved, never inferred just from compact policy sets. A selected-equilibrium target uses a declared selection rule; a robust target quantifies over all permitted equilibria. Use suprema or infima unless attainment is established. An \(\varepsilon\)-optimal policy is feasible and within \(\varepsilon\) of the finite optimum value. Report an empty feasible set separately.

\textbf{Welfare and accounting.} Normative utility, interpersonal normalization, social weights, numeraire, discounting and the use of public receipts are primitives. Behavioral utility need not coincide with the welfare criterion. Household welfare includes ownership income and rents once; adding profits again to compensating variation generally double-counts them. A monetary equivalent is defined by an explicit transfer and price convention, with monotonicity and range conditions. Average effects, mechanism contributions, transfers and welfare losses are different targets. Infinite sums require convergence and dynamic feasible plans require terminal or no-Ponzi conditions.

\textbf{Source versions.} A working-paper date, revision date and journal publication year are distinguished. A DOI for a journal version is not automatically the DOI of an earlier manuscript. Locators refer to the stated version and printed pages where specified. A metadata-only check does not verify a claim about a full-text conclusion. Every entry's source subsection narrows the provenance claim accordingly.
% END ECONOMICS STATEMENT
\end{document}
